% =========================================================
% Lab 2 — Decomposition (statement)
% =========================================================
\documentclass[a4paper,11pt]{article}

\newcommand{\TPnumero}{Lab 2 -- Decomposition}
% =========================================================
% Tutorial / lab preamble — Time Series Analysis module
% article class + fancyhdr + tcolorbox + listings (English)
% Author : Aymen Ben Brik (aymen.benbrik@esprit.tn)
% =========================================================

\usepackage[utf8]{inputenc}
\usepackage[T1]{fontenc}
\usepackage[english]{babel}
\usepackage{lmodern}
\usepackage{amsmath, amssymb, amsthm, mathtools, bm}
\usepackage{graphicx}
\usepackage{booktabs}
\usepackage{array}
\usepackage{multirow}
\usepackage{colortbl}
\usepackage{xcolor}
\usepackage{tikz}
\usetikzlibrary{positioning, arrows.meta, calc, shapes, fit, backgrounds, matrix}
\usepackage{listings}
\usepackage[most]{tcolorbox}
\usepackage{geometry}
\geometry{a4paper, margin=2cm, headheight=15pt}
\usepackage{fancyhdr}
\usepackage[hidelinks]{hyperref}
\usepackage{enumitem}

% --- Colours ---
\definecolor{espritBlue}{RGB}{30, 60, 110}
\definecolor{espritAccent}{RGB}{200, 80, 30}
\definecolor{boxEx}{RGB}{30, 60, 110}
\definecolor{boxQ}{RGB}{180, 120, 0}
\definecolor{boxC}{RGB}{30, 100, 60}
\definecolor{boxAtt}{RGB}{200, 80, 30}

% --- Listings ---
\definecolor{codebg}{RGB}{248,248,248}
\definecolor{codeframe}{RGB}{220,220,220}
\definecolor{keyword}{RGB}{0,82,155}
\definecolor{comment}{RGB}{88,110,117}
\definecolor{string}{RGB}{133,153,0}

\lstdefinestyle{pythonstyle}{
  language=Python,
  basicstyle=\ttfamily\small,
  keywordstyle=\color{keyword}\bfseries,
  commentstyle=\color{comment}\itshape,
  stringstyle=\color{string},
  backgroundcolor=\color{codebg},
  frame=single,
  rulecolor=\color{codeframe},
  frameround=tttt,
  breaklines=true,
  showstringspaces=false,
  tabsize=2
}
\lstdefinestyle{Rstyle}{
  language=R,
  basicstyle=\ttfamily\small,
  keywordstyle=\color{keyword}\bfseries,
  commentstyle=\color{comment}\itshape,
  stringstyle=\color{string},
  backgroundcolor=\color{codebg},
  frame=single,
  rulecolor=\color{codeframe},
  frameround=tttt,
  breaklines=true,
  showstringspaces=false,
  tabsize=2
}
\lstset{style=pythonstyle}

% --- Common macros (configurable path) ---
\providecommand{\macrosPath}{../../preamble/macros.tex}
\input{\macrosPath}

% --- Exercise counter ---
\newcounter{excounter}
\renewcommand{\theexcounter}{\arabic{excounter}}

\newtcolorbox[use counter=excounter]{exercise}[1][]{
  enhanced, breakable,
  colback=boxEx!5, colframe=boxEx, fonttitle=\bfseries,
  title={Exercise~\theexcounter\ifstrempty{#1}{}{ -- #1}},
  arc=2pt, boxrule=0.6pt
}
\newtcolorbox{question}[1][]{
  enhanced, breakable,
  colback=boxQ!5, colframe=boxQ, fonttitle=\bfseries,
  title={Question\ifstrempty{#1}{}{ -- #1}},
  arc=2pt, boxrule=0.5pt
}
\newtcolorbox{solution}[1][]{
  enhanced, breakable,
  colback=boxC!5, colframe=boxC, fonttitle=\bfseries,
  title={Solution\ifstrempty{#1}{}{ -- #1}},
  arc=2pt, boxrule=0.5pt
}
\newtcolorbox{warning}[1][]{
  enhanced, breakable,
  colback=boxAtt!5, colframe=boxAtt, fonttitle=\bfseries,
  title={Warning\ifstrempty{#1}{}{ -- #1}},
  arc=2pt, boxrule=0.5pt
}

% --- Headers / footers ---
\pagestyle{fancy}
\fancyhf{}
\fancyhead[L]{\textsc{\moduleNom} -- \auteurNom}
\fancyhead[R]{\TPnumero}
\fancyfoot[C]{\thepage}
\fancyfoot[L]{\auteurInstitution}
\fancyfoot[R]{\auteurEmail}
\renewcommand{\headrulewidth}{0.4pt}
\renewcommand{\footrulewidth}{0.2pt}

% --- Placeholder ---
\providecommand{\TPnumero}{Lab}


\title{Lab 2: Time series decomposition\\
\large Trend and seasonality on a synthetic series and on Atlanta temperatures}
\author{\auteurNom\\\auteurEmail\\\auteurInstitution}
\date{\anneeUniversitaire}

\begin{document}
\maketitle

\section*{Goals}
\begin{itemize}
  \item Implement and compare three trend estimators
        (centred moving average, polynomial regression, LOESS).
  \item Implement and compare three seasonality estimators
        (seasonal averaging, dummy regression, Fourier model).
  \item Apply the full workflow to the Atlanta monthly
        temperature dataset.
  \item Compare additive and multiplicative decomposition models.
\end{itemize}

\section*{Setup}
\begin{itemize}
  \item Python 3 with \texttt{numpy}, \texttt{pandas}, \texttt{scipy},
        \texttt{statsmodels}, \texttt{matplotlib}.
  \item Reproducibility: \texttt{numpy.random.seed(42)}.
  \item Dataset: \texttt{data/AvTempAtlanta.txt}.
  \item Skeleton: \texttt{seed.py} provides function signatures.
\end{itemize}

% =========================================================
\begin{exercise}[Synthetic decomposition]
Generate $T = 240$ months ($20$ years) of a synthetic series:
\[
  Y_t = \underbrace{20 + 0.05\, t}_{\text{trend}}
      + \underbrace{5 \sin(2\pi t / 12)}_{\text{seasonality}}
      + \underbrace{\varepsilon_t}_{\mathcal{N}(0, 1)}.
\]

\textbf{(1)} Plot $Y_t$ along with the true components on a 4-panel figure
(observed $Y_t$, true $T_t$, true $S_t$, true $\varepsilon_t$).

\textbf{(2)} Comment on the visual signature of the additive model
(constant seasonal amplitude).
\end{exercise}

% =========================================================
\begin{exercise}[Trend estimation: three methods]
\textbf{(1)} Implement the centred moving average
\texttt{centred\_ma(y, k)} of window $2k + 1$.

\textbf{(2)} Apply three methods to $Y_t$ from Exercise 1:
\begin{itemize}
  \item centred MA with $k = 6$ (window 13);
  \item linear regression (\texttt{sklearn.LinearRegression});
  \item LOESS (\texttt{statsmodels.nonparametric.smoothers\_lowess}, frac=0.2).
\end{itemize}
Plot the three estimated trends together with the true trend
$\widehat T_t = 20 + 0.05\,t$.

\textbf{(3)} Compute the RMSE of each estimator versus the true trend.
Comment.
\end{exercise}

% =========================================================
\begin{exercise}[Seasonality estimation: three methods]
\textbf{(1)} Detrend $Y_t$ from Exercise 1 using LOESS, then estimate
the seasonality by:
\begin{itemize}
  \item seasonal averaging (12 monthly indices);
  \item OLS regression with 11 dummy variables;
  \item cosine--sine model with one harmonic ($\omega = 2\pi / 12$).
\end{itemize}

\textbf{(2)} Plot the three estimated seasonal patterns
($\widehat S_k, k = 1, \dots, 12$) along with the truth
$5 \sin(2\pi k / 12)$.

\textbf{(3)} Compute the RMSE of each estimator versus the truth.
Which method performs best, and why?
\end{exercise}

% =========================================================
\begin{exercise}[Application to Atlanta temperatures]
Load \texttt{data/AvTempAtlanta.txt} (138 years $\times$ 12 months
= 1656 observations, time index $t = 1, \dots, 1656$).

\textbf{(1)} Plot the series. What is the seasonal period? Does the
amplitude look constant over time?

\textbf{(2)} Decompose using the workflow:
\begin{itemize}
  \item estimate the trend with a $2 \times 12$ centred MA;
  \item compute the detrended series;
  \item estimate seasonality by averaging (12 monthly indices, re-centred);
  \item compute the residual.
\end{itemize}
Plot the four components on a 4-panel figure.

\textbf{(3)} Compare with \texttt{statsmodels.tsa.seasonal\_decompose(y,
period=12, model="additive")}. Do you get the same components?
\end{exercise}

% =========================================================
\begin{exercise}[Bonus: additive vs multiplicative]
\textbf{(1)} Generate a multiplicative series:
\[
  Y_t = (20 + 0.1\, t) \cdot \bigl(1 + 0.3 \sin(2\pi t / 12)\bigr) \cdot e^{\eta_t},
  \quad \eta_t \sim \mathcal{N}(0, 0.05).
\]
Plot it. Comment on the seasonal amplitude.

\textbf{(2)} Apply the additive workflow on $Y_t$ \emph{vs}.\ the
additive workflow on $\log Y_t$. Compare the residuals.

\textbf{(3)} Conclude: when should you use the multiplicative model?
\end{exercise}

\bigskip
\begin{warning}
For each plot, save a PNG (\texttt{lab2\_ex1\_synth.png}, etc.). Set
the seed at the top of every script.
\end{warning}

\end{document}
